\newcommand{\figuraCeldaPresion}{%
\begin{figure}[htbp]
    \centering
    \begin{tikzpicture}[>=Latex,font=\small]
        \coordinate (A) at (2.10,0.85);
        \coordinate (B) at (6.65,0.85);
        \coordinate (C) at (6.65,4.45);
        \coordinate (D) at (2.10,4.45);
        \coordinate (E) at (3.35,1.80);
        \coordinate (F) at (7.90,1.80);
        \coordinate (G) at (7.90,5.40);
        \coordinate (H) at (3.35,5.40);

        \fill[blue!5] (A)--(B)--(C)--(D)--cycle;
        \fill[blue!12] (B)--(F)--(G)--(C)--cycle;
        \fill[blue!18] (D)--(C)--(G)--(H)--cycle;
        \draw[very thick] (A)--(B)--(C)--(D)--cycle;
        \draw[very thick] (B)--(F)--(G)--(C);
        \draw[very thick] (C)--(G)--(H)--(D);
        \draw[very thick,dashed] (A)--(E)--(F);
        \draw[very thick,dashed] (E)--(H);

        \node[font=\bfseries] at (5.00,6.42)
            {Presión sobre una celda cúbica};
        \node[text=black!70] at (4.38,2.55)
            {$p=p(\vec{r},t)$};

        % Cara superior: vector de superficie y fuerza de presión.
        \coordinate (T) at (5.18,4.92);
        \draw[very thick,blue!70!black,-{Latex[length=3mm]}]
            (T)--(5.18,5.95);
        \node[blue!70!black,anchor=west] at (5.28,5.66)
            {$d\vec{S}_z$};
        \draw[very thick,red!75!black,-{Latex[length=3mm]}]
            (4.72,5.92)--(4.72,5.03);
        \node[red!75!black,anchor=east] at (4.62,5.60)
            {$d\vec{F}_{p,z}$};

        % Cara derecha.
        \coordinate (R) at (7.28,3.30);
        \draw[very thick,blue!70!black,-{Latex[length=3mm]}]
            (R)--(8.73,3.72);
        \node[blue!70!black,anchor=west] at (8.73,3.72)
            {$d\vec{S}_x$};
        \draw[very thick,red!75!black,-{Latex[length=3mm]}]
            (8.60,3.16)--(7.42,2.82);
        \node[red!75!black,anchor=west] at (8.52,2.88)
            {$d\vec{F}_{p,x}$};

        % Dimensiones de la celda.
        \draw[<->,thick] (2.10,0.48)--node[below]{$dx$}(6.65,0.48);
        \draw[<->,thick] (1.72,0.85)--node[left]{$dz$}(1.72,4.45);
        \draw[<->,thick] (6.87,0.65)--node[below,sloped]{$dy$}(8.12,1.60);

        \node[align=left,anchor=west] at (9.20,4.90)
            {\textcolor{blue!70!black}{\rule{0.55cm}{1.2pt}}\;
             vector de superficie\\[5pt]
             \textcolor{red!75!black}{\rule{0.55cm}{1.2pt}}\;
             fuerza de presión};
        \node[align=center,anchor=west] at (9.20,2.48)
            {$d\vec{F}_p=-p\,d\vec{S}$\\[4pt]
             La fuerza apunta\\hacia el interior};
    \end{tikzpicture}
    \caption{Vectores de superficie y fuerzas de presión sobre dos caras de
    una celda cúbica. Cada vector $d\vec{S}$ apunta hacia afuera, mientras que
    la fuerza $d\vec{F}_p=-p\,d\vec{S}$ actúa normal a la cara y hacia el
    interior. La presión isotrópica tiene el mismo valor para todas las
    orientaciones.}
    \label{fig:celda-presion}
\end{figure}%
}
